% ============================================================
%  Kavya Prakash — Industry Application CV
%  Entry-level | EU Life Sciences (Biotech / Pharma / CRO / Genomics)
%  ATS-safe, single-column, clean minimal style
%  Compiler: pdfLaTeX  |  Overleaf ready
% ============================================================

\documentclass[10pt, a4paper]{article}

% ── Packages ────────────────────────────────────────────────
\usepackage[
  top=1.8cm, bottom=1.8cm,
  left=2cm,  right=2cm
]{geometry}
\usepackage{enumitem}
\usepackage{titlesec}
\usepackage{parskip}
\usepackage{array}
\usepackage{tabularx}
\usepackage{microtype}
\usepackage[hidelinks]{hyperref}
\usepackage[T1]{fontenc}
\usepackage[utf8]{inputenc}
\usepackage{lmodern}

% ── Typography ───────────────────────────────────────────────
\pagestyle{empty}
\setlength{\parindent}{0pt}
\setlength{\parskip}{0pt}
\hypersetup{urlcolor=black, linkcolor=black}

% Section headers: bold, all-caps, ruled line underneath
\titleformat{\section}
  {\large\bfseries\uppercase}
  {}{0em}{}
  [\vspace{1pt}\hrule\vspace{4pt}]
\titlespacing{\section}{0pt}{10pt}{4pt}

% Bullet list settings
\setlist[itemize]{
  leftmargin=1.4em,
  itemsep=1pt,
  parsep=0pt,
  topsep=2pt
}

% ── Helper commands ──────────────────────────────────────────
\newcommand{\entry}[3]{%
  \begin{tabularx}{\linewidth}{@{}X r@{}}
    \textbf{#1} & \textit{#3} \\
    \textit{#2} & \\
  \end{tabularx}%
}

\newcommand{\entryflat}[2]{%
  \begin{tabularx}{\linewidth}{@{}X r@{}}
    \textbf{#1} & \textit{#2} \\
  \end{tabularx}%
}

% ============================================================
\begin{document}
% ============================================================

% ── Header ──────────────────────────────────────────────────
\begin{center}
  {\LARGE\bfseries Kavya Prakash}\\[4pt]
  Uppsala, Sweden \quad|\quad
  +46 7696 40 800 \quad|\quad
  \href{mailto:kavyaprakash0398@gmail.com}{kavyaprakash0398@gmail.com}\\[2pt]
  \href{https://kavyaprakash.me}{Portfolio}
  \quad|\quad
  \href{https://www.linkedin.com/in/kavya-prakash-vks/}{LinkedIn}
  \quad|\quad
  \href{https://github.com/kavyaprakash-vks}{GitHub}
\end{center}

\vspace{2pt}

% ── Professional Profile ─────────────────────────────────────
\section{Professional Profile}

M.Sc. graduate in Biology (Immunology and Microbiology) from Uppsala University with hands-on
experience in stem cell culture, molecular assay development, and transcriptomic data analysis.
Proven ability to design, optimise, and execute reproducible wet-lab workflows across neuronal
differentiation, epigenetics, and functional genomics. Adept at working in collaborative,
multi-project research environments and communicating scientific findings to mixed audiences.
Seeking an entry-level Research Associate or Scientist position in the EU life sciences sector.

\vspace{2pt}

% ── Technical Skills ─────────────────────────────────────────
\section{Technical Skills}

\begin{tabularx}{\linewidth}{@{} >{\bfseries}p{5cm} X @{}}
Molecular Assays
  & qPCR, RT-qPCR, pyrosequencing (bisulfite-based DNA methylation), RNA extraction,
    DNA extraction, immunofluorescence staining, ELISA, Western blotting \\[3pt]

Cell Culture \& Models
  & Human and mouse embryonic stem cells (hESC, mESC), SH-SY5Y neuroblastoma cells,
    Caco-2 intestinal epithelial cells, dorsal forebrain differentiation protocols \\[3pt]

Additional Techniques
  & Flow cytometry, whole-mount in situ hybridisation, live-cell imaging,
    GLP-compliant practice, QA/QC procedures \\[3pt]

Bioinformatics \& Data
  & R (DESeq2, ggplot2), RNA-seq analysis, differential gene expression,
    Gene Ontology and KEGG enrichment analysis, PCA, heatmaps, Linux command line \\[3pt]

Laboratory Operations
  & Autoclaving, glassware washing, waste disposal, cell culture hood and equipment
    upkeep, consumables restocking, and supply ordering/vendor communication \\[3pt]

Languages
  & English (IELTS 8.0), Malayalam (native), Swedish (basic) \\
\end{tabularx}

\vspace{2pt}

% ── Scientific Experience ────────────────────────────────────
\section{Scientific Experience}

\entry{Master's Thesis Researcher}{EpiTox Group, Department of Organismal Biology, Uppsala University}{Oct 2025 -- May 2026}
\vspace{2pt}
\textit{Cellular and Molecular Impact of BPA and BPF on Human and Mouse Embryonic Stem Cell Models of Dorsal Forebrain Differentiation}
\begin{itemize}
  \item Established and optimised a dorsal forebrain differentiation protocol in hESC and mESC
        to model early cortical neurodevelopment in vitro, enabling reproducible cell-based
        assay development.
  \item Characterised the effects of endocrine-disrupting chemicals (BPA, BPF) on
        neuroectodermal differentiation markers using RT-qPCR and immunofluorescence staining.
  \item Integrated multi-modal molecular and imaging datasets to evaluate EDC-mediated
        disruption of epigenetic and transcriptional programmes during neural induction.
\end{itemize}

\vspace{5pt}

\entry{Laboratory Teaching Assistant}{Uppsala University}{Nov 2025 -- Jan 2026}
\begin{itemize}
  \item Trained undergraduate and graduate students in pyrosequencing, RT-qPCR, zebrafish
        whole-mount in situ hybridisation, and live-cell imaging; translated complex protocols
        into accessible, step-by-step practical guidance.
  \item Managed sample preparation, equipment setup, and experimental workflow logistics for
        multi-student laboratory sessions.
\end{itemize}

\vspace{5pt}

\entry{Research Trainee}{Department of Organismal Biology --- Physiology and Environmental Toxicology, Uppsala University}{Jul 2025 -- Sep 2025}
\begin{itemize}
  \item Differentiated SH-SY5Y neuroblastoma cells toward a glutamatergic neuronal phenotype
        and validated marker expression using qPCR; results informed downstream toxicological
        assay design.
  \item Maintained SH-SY5Y and Caco-2 cell lines and optimised molecular protocols to support
        concurrent studies across multiple research projects.
  \item Managed laboratory maintenance and operations including autoclaving, glassware washing,
        waste disposal, cell culture equipment upkeep, and consumables restocking.
\end{itemize}

\vspace{5pt}

\entry{Quality Assurance \& Quality Control Specialist}{Nexus Agro Farms, Kerala, India}{Oct 2021 -- Jun 2024}
\begin{itemize}
  \item Applied QA/QC procedures and data-driven analysis to support product consistency and
        regulatory compliance across testing operations.
  \item Maintained testing records, batch documentation, and compliance reports; coordinated
        with suppliers and internal teams to resolve quality deviations and ensure traceability.
\end{itemize}

\vspace{2pt}

% ── Independent Research Project ─────────────────────────────
\section{Independent Research Project}

\entryflat
{Transcriptomic Analysis of Valproic Acid-Induced Neurodevelopmental Alterations in Human Forebrain Organoids}
{Independent Bioinformatics Project}

\vspace{3pt}

\begin{itemize}
\item Re-analysed publicly available RNA-seq data from human forebrain organoids exposed to
      valproic acid (VPA) using a reproducible end-to-end bioinformatics pipeline.
\item Performed transcript quantification (Salmon), differential expression analysis (DESeq2),
      and functional enrichment (GO, KEGG) entirely in R; identified 7,151 significantly
      differentially expressed genes (FDR $<$ 0.05), of which 1,676 showed absolute
      log$_2$ fold change $>$ 1.
\item Repository: \url{https://github.com/kavyaprakash-vks/vpa-forebrain-organoid-rnaseq}
\end{itemize}



\vspace{2pt}

% ── Education ────────────────────────────────────────────────
\section{Education}

\entry{M.Sc.\ in Biology --- Immunology and Microbiology}{Uppsala University, Sweden}{2024 -- 2026}

\vspace{4pt}

\entry{M.Sc.\ in Biotechnology}{Kerala University, India}{2019 -- 2021}
\vspace{1pt}
\textit{Thesis: Identification of microRNAs Regulating Differentially Expressed Genes in Cisplatin-Resistant Ovarian Cancer Cells}
\begin{itemize}
  \item Profiled microRNA-mediated regulation of chemoresistance pathways in ovarian cancer
        using expression analysis and bioinformatics-assisted pathway interpretation.
  \item Identified candidate miRNA--mRNA regulatory axes associated with platinum-based drug
        resistance mechanisms.
\end{itemize}

\vspace{4pt}

\entry{B.Sc.\ in Biotechnology}{Kerala University, India}{2016 -- 2019}

\vspace{2pt}

% ── Additional ───────────────────────────────────────────────
\section{Additional}

\entryflat{International Student Ambassador, Office of Science and Technology --- Uppsala University}{2025 -- 2026}

\vspace{4pt}

\textbf{Conference Presentation:} Kavya Prakash. ``Environmental Pollutants and Brain Health:
Insights from Human Cell Models of Neurodegeneration.''
\textit{Uppsala University Student Conference in Science and Technology}, 2025.

\vspace{2pt}

% ── References ───────────────────────────────────────────────
\section{References}

Available upon request.

% ============================================================
\end{document}
